% Part: lambda-calculus
% Chapter: syntax
% Section: abbreviated-syntax

\documentclass[../../../include/open-logic-section]{subfiles}

\begin{document}

\olfileid{lam}{syn}{abb}
\olsection{نحو اختصاری}

نوشتن ترم‌هایی که در~\olref[trm]{defn:term} تعریف شده‌اند گاهی
دست‌وپاگیر است، پس معرفی نحوی موجزتر سودمند است. البته باید مراقب باشیم
که ترم‌ها در نمادگذاری موجز نیز خوانش یکتا داشته باشند. یکی از نسخه‌های
پرکاربرد که \emph{ترم‌های اختصاری} نامیده می‌شود چنین است.

\begin{enumerate}
\item وقتی پرانتزها حذف شوند، کاربرد از چپ به راست انجام می‌شود. برای
  مثال، اگر~$M$، $N$، $P$ و~$Q$ ترم باشند، $MNPQ$ اختصار
  $(((MN)P)Q)$ است.
\item باز هم وقتی پرانتزها حذف شوند، انتزاع لامبدایی گسترده‌ترین دامنهٔ
  ممکن را می‌گیرد. برای مثال، $\lambd[x][MNP]$ به‌صورت
  $(\lambd[x][MNP])$ خوانده می‌شود.
\item می‌توان از یک لامبدا برای انتزاع چند متغیر استفاده کرد. برای مثال،
  $\lambd[xyz][M]$ اختصار
  $\lambd[x][\lambd[y][\lambd[z][M]]]$ است.
\end{enumerate}

برای مثال،
\[
\lambd[xy][xxyx \lambd[z][xz]]
\]
اختصار عبارت زیر است:
\[
(\lambd[x][(\lambd[y][((((xx)y)x)(\lambd[z][(xz)]))])]).
\]

\begin{prob}
ترم اختصاری~$\lambd[g][(\lambd[x][g (x x)])
  \lambd[x][g (x x)]]$ را بسط دهید.
\end{prob}

\end{document}
